% Alpha_Number_Theory.tex
% Deep investigation: number-theoretic properties of the DFD alpha formula
% Date: 2026-03-23

\documentclass[11pt]{article}
\usepackage{amsmath,amssymb,amsthm}
\usepackage[margin=1in]{geometry}
\usepackage{booktabs,array}

\newcommand{\Z}{\mathbb{Z}}
\newcommand{\Q}{\mathbb{Q}}

\title{Number-Theoretic Properties of the DFD\\
Fine Structure Constant Formula}
\author{DFD Research Programme}
\date{23 March 2026}

\begin{document}
\maketitle

%% =========================================================================
\section{The Formula and Its Integers}
\label{sec:formula}
%% =========================================================================

The DFD spectral action pipeline yields (with $\mathrm{Tr}(Y^2)=10$, $k=60$):
\begin{equation}
\alpha^{-1}
\;=\;
\frac{\pi^{3/2}}{24}\times 10 \times 60
\times \frac{63}{64}
\times \biggl(1 + \frac{7}{327600}\biggr)
\;=\; 137.035\,999\,854\ldots
\label{eq:master}
\end{equation}
which agrees with the CODATA 2018 value $137.035\,999\,084(21)$ to 5.6~ppb.

The integers appearing are: $24,\;10,\;60,\;63,\;64,\;7,\;327600$.

%% =========================================================================
\section{Prime Factorizations}
\label{sec:primes}
%% =========================================================================

\begin{center}
\renewcommand{\arraystretch}{1.3}
\begin{tabular}{r@{\;=\;}l@{\qquad}l}
\toprule
\textbf{Integer} & \textbf{Factorization} & \textbf{Role in formula} \\
\midrule
$24$     & $2^3\times 3 = 4!$              & $1/a_4$ normalisation \\
$10$     & $2\times 5$                      & $\mathrm{Tr}(Y^2)$, SM hypercharges \\
$60$     & $2^2\times 3\times 5$            & Spectral cutoff $k_{\max}$ \\
$63$     & $3^2\times 7 = 2^6-1$            & Numerator of $(k+3)/(k+4)$ \\
$64$     & $2^6$                            & Denominator $k+4 = 64$ \\
$7$      & $7$ (prime)                      & Numerator of radiative correction \\
$327600$ & $2^4\times 3^2\times 5^2\times 7\times 13$ & Denominator of correction \\
\bottomrule
\end{tabular}
\end{center}

\medskip\noindent
\textbf{Observation.} The only primes that enter the \emph{entire} formula are
$\{2,3,5,7,13\}$---the first five primes \emph{except}~11.
The prime~13 enters solely through $327600$ (specifically through $k+4=64$,
since $(k+4)^2-1 = 4095 = 3^2\times 5\times 7\times 13$).

%% =========================================================================
\section{Factorial and Mersenne Structure}
\label{sec:factorial}
%% =========================================================================

\subsection{The prefactor: $(5!)^2/4!$}

The combination $10\times 60/24 = 25$ can be rewritten using factorials:
\begin{equation}
\frac{10\times 60}{24}
= \frac{600}{4!}
= \frac{(5!)^2}{(4!)^2}\cdot\frac{4!}{(5!)^{-1}\cdots}
= 5^2 = 25.
\end{equation}
More revealingly:
\[
24\times 10\times 60 = 14400 = 120^2 = (5!)^2.
\]
The product of the three base integers is the square of $5! = |S_5|$.

\subsection{The $2^6$ architecture}

Both correction terms involve $2^6 = 64$:
\begin{align}
\frac{63}{64} &= \frac{2^6 - 1}{2^6} = 1 - 2^{-6},
\label{eq:mersenne6}\\[4pt]
\frac{7}{327600}
&= \frac{7}{7!\times 65}
= \frac{1}{6!\times(2^6+1)}.
\label{eq:correction}
\end{align}

\noindent
\textbf{Key identity:} $327600 = 7!\times 65 = 7!\times(2^6+1)$,
and $4095 = 63\times 65 = (2^6-1)(2^6+1) = 2^{12}-1$.

This means the formula's two sub-leading corrections are:
\begin{itemize}
\item A \textbf{Mersenne factor} $(2^6-1)/2^6$, involving the 6th Mersenne number.
\item A \textbf{fine correction} $1/(6!\times(2^6+1))$, pairing a factorial with
      the conjugate Mersenne cofactor.
\end{itemize}
Together they produce the 12-bit Mersenne number $2^{12}-1=4095$ in the denominator.

\subsection{Factorial form}

Combining the above:
\begin{equation}
\boxed{\;
\alpha^{-1}
= \frac{\pi^{3/2}\times(5!)^2}{(4!)^2}
\;\times\;\frac{2^6-1}{2^6}
\;\times\;\biggl(1 + \frac{1}{6!\,(2^6+1)}\biggr)
\;}
\label{eq:factorial}
\end{equation}

\subsection{Simplest rational form}

Expanding all rational factors and cancelling:
\begin{equation}
\alpha^{-1}
= \frac{\pi^{3/2}\times 327607}{2^{10}\times 13}
= \frac{327607\,\pi^{3/2}}{13312},
\label{eq:simplest}
\end{equation}
where $327607 = 7\times 17\times 2753$ (with 2753 prime),
and $13312 = 2^{10}\times 13$.

%% =========================================================================
\section{Connection to Bernoulli Numbers}
\label{sec:bernoulli}
%% =========================================================================

By the von~Staudt--Clausen theorem, the denominator of the $n$-th Bernoulli
number $B_n$ is $\prod_{(p-1)|n} p$.  For $B_{12}$:
\[
(p-1)\mid 12 \implies p\in\{2,3,5,7,13\}
\implies \mathrm{denom}(B_{12}) = 2\times 3\times 5\times 7\times 13 = 2730.
\]

\noindent
\textbf{Discovery:}
\begin{equation}
327600 \;=\; \mathrm{denom}(B_{12})\times 5! \;=\; 2730\times 120.
\label{eq:bernoulli}
\end{equation}
The set of primes $\{2,3,5,7,13\}$ dividing $327600$ is \emph{exactly}
the set of primes appearing in $\mathrm{denom}(B_{12})$.

The Bernoulli numbers $B_{2n}$ control the constant terms of Seeley--DeWitt
(heat-kernel) coefficients $a_{2n}$, which are precisely the objects integrated
in the DFD spectral action.  That the correction denominator factorises as
$\mathrm{denom}(B_{12})\times 5!$ may therefore not be an accident: the
spectral action at order $a_4$ naturally generates rational prefactors whose
prime support is inherited from $B_{12}$ via the heat-kernel expansion.

Additionally, $4095 = (3/2)\times\mathrm{denom}(B_{12})$, tightening the link
between the Mersenne structure and the Bernoulli denominators.

%% =========================================================================
\section{Group-Theoretic Identifications}
\label{sec:groups}
%% =========================================================================

\begin{center}
\renewcommand{\arraystretch}{1.3}
\begin{tabular}{r@{\;=\;}ll}
\toprule
\textbf{Integer} & \textbf{Group order} & \textbf{Group} \\
\midrule
$24$  & $|S_4|$     & Symmetric group on 4 letters; chiral octahedral group \\
$60$  & $|A_5|$     & Alternating group; icosahedral rotation group; $\mathrm{PSL}(2,5)$ \\
$120$ & $|S_5|$     & Full icosahedral group (with reflections); $\mathrm{SL}(2,5)$ \\
$240$ & $|$roots of $E_8|$ & Root system of $E_8$ ($24\times 10 = 240$) \\
\bottomrule
\end{tabular}
\end{center}

\medskip\noindent
\textbf{Is $60 = |A_5|$ a coincidence?}  The value $k_{\max}=60$ is the spectral
truncation of the DFD action.  The alternating group $A_5$ is the \emph{unique}
finite simple group of order~60, and is isomorphic to the icosahedral rotation
group.  Whether the spectral cutoff is forced to the order of the icosahedral
group by some deeper geometric reason (e.g., the icosahedral subgroup of $\mathrm{SO}(3)$
controlling angular-momentum truncation on $S^3$) is an open question.

\textbf{E${}_8$ roots.}  The product $24\times 10 = 240$ equals the number of
roots of the $E_8$ lattice.  In the DFD formula this corresponds to
$a_4\text{-normalisation}\times\mathrm{Tr}(Y^2)$. If the hypercharge trace
$\mathrm{Tr}(Y^2)=10$ is ultimately explained by an $E_8$ embedding of the
Standard Model gauge group, this coincidence becomes structural.

%% =========================================================================
\section{Properties of 137}
\label{sec:137}
%% =========================================================================

\begin{itemize}
\item \textbf{33rd prime.}  $137 = p_{33}$, where 33 = $3\times 11$.
\item \textbf{Pythagorean prime.}  $137 \equiv 1\pmod{4}$, so by Fermat's theorem
      $137 = 4^2 + 11^2$.
\item \textbf{Binary.}  $137 = (10001001)_2 = 2^7 + 2^3 + 2^0$.
      Hamming weight~3; bit positions $\{0,3,7\}$.
\item \textbf{Digit sum.}  $1+3+7 = 11$ (prime).
\item \textbf{Truncatable prime.}  137 is both left-truncatable (137, 37, 7 all prime)
      and right-truncatable (137, 13, 1$\ldots$ well, 1 is not prime; it is
      left-truncatable only).
\item \textbf{Golden angle.}  $\lfloor 360/\varphi^2\rfloor = \lfloor 137.508\ldots\rfloor = 137$,
      where $\varphi = (1+\sqrt{5})/2$.
\item \textbf{Palindromic reciprocal.}  $1/137 = 0.\overline{00729927}$ with
      period~8.  The repeating block is a palindrome.
\end{itemize}

\noindent
\textbf{Connection to the formula integers?}
The continued fraction of $\alpha^{-1}$ is $[137;\,27,\,1,\,3,\,1,\,1,\,16,\,1,\,10,\,3,\ldots]$.
The second partial quotient, 27, yields the excellent rational approximant
$3700/27 = 137.037\overline{037}$, accurate to $10^{-3}$.
The convergent $34259/250 = 137.036\,000\,000$ is exact to $10^{-6}$.
There is no obvious direct arithmetic link between 137 and the formula's
integers $\{24,10,60,63,64\}$, consistent with 137 being an \emph{emergent}
consequence of irrational $\pi^{3/2}$ multiplied by the rational prefactor,
rather than a built-in input.

%% =========================================================================
\section{Continued Fraction of $\alpha^{-1}$}
\label{sec:cf}
%% =========================================================================

\begin{equation}
\alpha^{-1} = [137;\,27,\,1,\,3,\,1,\,1,\,16,\,1,\,10,\,3,\,1,\,2,\,1,\,4,\,\ldots]
\end{equation}

\noindent
Key convergents:
\begin{center}
\begin{tabular}{crl}
\toprule
\textbf{Order} & \textbf{Convergent} & \textbf{Error} \\
\midrule
0 & $137/1$          & $-3.6\times 10^{-2}$ \\
1 & $3700/27$        & $+1.0\times 10^{-3}$ \\
5 & $34259/250$      & $+9.2\times 10^{-7}$ \\
6 & $567192/4139$    & $-5.0\times 10^{-8}$ \\
8 & $6581702/48029$  & $-1.1\times 10^{-10}$ \\
\bottomrule
\end{tabular}
\end{center}

\noindent
The continued fraction is unremarkable---no exceptionally large partial
quotients, no periodic structure.  This confirms that $\alpha^{-1}$ is not a
quadratic irrational.  Given the DFD formula~\eqref{eq:simplest}, $\alpha^{-1}$
is a rational multiple of $\pi^{3/2}$; its continued fraction inherits
the ``generic'' behaviour of $\pi^{3/2}$, which is expected to be normal
(all partial quotients appearing with Gauss--Kuzmin frequencies).

%% =========================================================================
\section{Is $25\pi^{3/2}$ Close to an Algebraic Number?}
\label{sec:algebraic}
%% =========================================================================

\[
25\pi^{3/2} = 139.208\,199\,920\,793\ldots
\]
Its continued fraction is $[139;\,4,\,1,\,4,\,12,\,1,\,4,\,8,\,5,\,2,\,\ldots]$,
again with no striking pattern.

\medskip\noindent
\textbf{Simple-fraction search.}
Scanning all fractions $p/q$ with $1\le q\le 100$:
the closest are $139/1$ (error $+0.208$), $557/4$ (error $-0.042$),
and $696/5$ (error $+0.792/5 = -0.008$).
None is remarkably close.

\medskip\noindent
\textbf{Conclusion.}
By the Lindemann--Weierstrass theorem, $\pi^{3/2}$ is transcendental,
so $25\pi^{3/2}$ is transcendental and cannot equal any algebraic number.
The question is whether it is \emph{unusually close} to one.  No such
proximity has been found; the quantity behaves as a ``generic'' transcendental.

%% =========================================================================
\section{Correction Hierarchy}
\label{sec:hierarchy}
%% =========================================================================

The formula has a natural three-layer structure:

\begin{center}
\begin{tabular}{llrl}
\toprule
\textbf{Level} & \textbf{Expression} & \textbf{Value} & \textbf{Residual} \\
\midrule
0 & $25\pi^{3/2}$                    & $139.208\ldots$ & $+1.59\%$ \\
1 & $\times\,(2^6-1)/2^6$            & $137.033\ldots$ & $-0.0021\%$ \\
2 & $\times\,(1+1/(6!\cdot 65))$     & $137.036\ldots$ & $5.6$~ppb \\
\bottomrule
\end{tabular}
\end{center}

\noindent
Each correction reduces the residual by a factor of $\sim\!700$--$400$.
The Level-1 correction $(2^6-1)/2^6$ removes $\sim\!1.6\%$;
the Level-2 correction $1/(6!\cdot 65) \approx 21$~ppm removes all but 5.6~ppb.

%% =========================================================================
\section{Summary of Number-Theoretic Content}
\label{sec:summary}
%% =========================================================================

\begin{enumerate}
\item \textbf{Primes used:} $\{2,3,5,7,13\}$ only (notably, 11 is absent).
      These are exactly the primes dividing $\mathrm{denom}(B_{12})$.

\item \textbf{Factorial skeleton:}
      $24 = 4!$,\; $(5!)^2 = 14400 = 24\times 10\times 60$,\;
      $7! = 5040$,\; $6! = 720$.
      The formula is built from $4!$, $5!$, $6!$, and $7!$.

\item \textbf{Mersenne/power-of-two skeleton:}
      $63 = 2^6-1$,\; $64 = 2^6$,\; $65 = 2^6+1$,\;
      $4095 = 2^{12}-1$,\; $13312 = 2^{10}\times 13$.
      The number~6 (from $k+4 = 64 = 2^6$) controls both corrections.

\item \textbf{Bernoulli connection:}
      $327600 = \mathrm{denom}(B_{12})\times 5!$, linking the correction
      denominator to the 12th Bernoulli number---natural for a heat-kernel
      (spectral action) calculation.

\item \textbf{Group theory:}
      $60 = |A_5|$ (icosahedral), $24 = |S_4|$ (octahedral),
      $240 = 24\times 10 = |$roots of $E_8|$.
      The product $(5!)^2$ is the square of $|S_5|$.

\item \textbf{137 itself} is an emergent output---the 33rd prime, a Pythagorean
      prime ($4^2+11^2$), and the integer part of the golden angle---but bears
      no direct arithmetic relation to the formula's input integers.

\item \textbf{Simplest closed form:}
      $\alpha^{-1} = 327607\,\pi^{3/2}/(2^{10}\times 13)$,
      with $327607 = 7\times 17\times 2753$.
\end{enumerate}

\end{document}
